\section*{Exercises about Polynomials}

\exercise{3}
\begin{xrcs}
  (Lagrange Interpolation) Suppose $z_1, \ddd, z_{m+1} \in \myF$ are distinct and $w_1, \ddd, w_{m+1} \in \myF$. Prove that there exists a unique polynomial $p \in \polyn_m(\myF)$ such that
  \begin{equation}
    p(z_j) = w_j \quad \text{for } j = 1, \ddd, m+1.
  \end{equation}
% \begin{xprf}
%   Define the linear evaluation map $T: \polyn_m(\myF) \to \myF^{m+1}$ by
%   \[
%     T(p) := (p(z_1), p(z_2), \dots, p(z_{m+1})).
%   \]
%   If $p \in \mynull T$, then $p(z_j) = 0$ for each $j \in \{1, \dots, m+1\}$.
%   Thus $p$ has at least $m+1$ distinct roots in $\myF$.
%   Since $\deg p \le m$, the only polynomial with more roots than its degree is the zero polynomial, so $p = 0$.
%   Thus $\mynull T = \{0\}$, meaning $T$ is injective.
%   Since $\dim \polyn_m(\myF) = m+1 = \dim \myF^{m+1}$, Theorem \ref{thm: injectivity is equivalent to surjectivity} implies that $T$ is an isomorphism (bijective).
%   Therefore, for every $(w_1, \dots, w_{m+1}) \in \myF^{m+1}$, there exists a unique $p \in \polyn_m(\myF)$ such that $T(p) = (w_1, \dots, w_{m+1})$, which proves $p(z_j) = w_j$ for each $j$.
% \end{xprf}
\end{xrcs}

\exercise{4}
\begin{xrcs}
  Suppose $m$ is a positive integer. Is the set
  \begin{equation}
    S := \{0\} \cup \{ p \in \polyn (\myF) : \deg p = m \}
  \end{equation}
  a subspace of $\polyn (\myF)$?

  \begin{xsol}
    The set $S$ is not a subspace of $\polyn (\myF)$ because it is not closed under addition.
    For example, let $p(x) = x^m + x$ and $q(x) = -x^m$. Both $p$ and $q$ have degree $m$ (so $p, q \in S$), but
    \[
      p(x) + q(x) = x
    \]
    has degree $1 \neq m$ (for $m \neq 1$) and is not $0$, so $p+q \notin S$. (If $m=1$, take $p(x)=x+1$ and $q(x)=-x$, so $p(x)+q(x)=1$, which has degree $0 \neq 1$). Thus, $S$ is not a subspace.
  \end{xsol}
\end{xrcs}

\exercise{5}
\begin{xrcs}
  Is the set
  \begin{equation}
    S := \{0\} \cup \{ p \in \polyn (\myF) : \deg p \text{ is even} \}
  \end{equation}
  a subspace of $\polyn (\myF)$?

  \begin{xsol}
    The set $S$ is not a subspace of $\polyn (\myF)$ because it is not closed under addition.
    For example, let $p(x) = x^4 + x^3$ and $q(x) = -x^4$. Both $p$ and $q$ have degree $4$ (which is even, so $p, q \in S$), but
    \[
      p(x) + q(x) = x^3
    \]
    has degree $3$, which is odd and not $0$. Thus, $p+q \notin S$, so $S$ is not a subspace.
  \end{xsol}
\end{xrcs}

\exercise{6}
\begin{xrcs}
  Suppose that $m$ and $n$ are positive integers with $m \leq n$, and suppose $\lambda_1, \ddd, \lambda_m \in \myF$. Prove that there exists a polynomial $p \in \polyn(\myF)$ with $\deg p = n$ such that $0 = p(\lambda_1) = \cdots = p(\lambda_m)$ and such that $p$ has no other zeros in $\myF$.

  \begin{xprf}
    Define $p \in \polyn(\myF)$ by
    \begin{equation}
      \begin{aligned}
        p(z) &:= (z-\lambda_1)(z-\lambda_2)\cdots(z-\lambda_{m-1}) \underbrace{(z-\lambda_m)(z-\lambda_m)\cdots(z-\lambda_m)}_{n-m+1 \text{ times}} \\
             &= \left(\prod_{k = 1}^{m-1} (z-\lambda_k) \right) (z-\lambda_m)^{n-m+1} \qquad (\forall z \in \myF).
      \end{aligned}
    \end{equation}
    The degree of $p$ is $(m-1) + (n-m+1) = n$.
    Clearly, evaluating $p$ at $z = \lambda_k$ yields $p(\lambda_k) = 0$ for every $k \in \{1, \ddd, m\}$.

    Now, suppose $\lambda \in \myF$ is a zero of $p$, so $p(\lambda) = 0$. Since $\myF$ is a field, the product
    \[
      (\lambda-\lambda_1)(\lambda-\lambda_2)\cdots(\lambda-\lambda_{m-1})(\lambda-\lambda_m)^{n-m+1} = 0
    \]
    can only equal $0$ if at least one of the factors is $0$.
    Thus, $\lambda - \lambda_k = 0$ for some $k \in \{1, \ddd, m\}$, meaning $\lambda = \lambda_k$.
    Therefore, $p$ has no zeros other than $\lambda_1, \ddd, \lambda_m$.
  \end{xprf}
\end{xrcs}

% ==============================================================================
% ORIGINAL ATTEMPT (Commented out): Exercise 6 Draft
% ------------------------------------------------------------------------------
% AI FEEDBACK:
% 1. Product Notation: Line 41 of the draft wrote $\sum_{k=1}^{m-1} (z-\lambda_k)$
%    instead of the product $\prod_{k=1}^{m-1} (z-\lambda_k)$.
% 2. Statement Consistency: The problem hypothesis requires $m \le n$ (and $\deg p = n$),
%    whereas the draft had $m \neq n$ in the premise.
% 3. Environment Scoping: The second part of the argument was placed outside \end{xprf}.
% ==============================================================================
% \begin{xprf}
%   Define $p \in \polyn_n (\myF)$ by
%   \begin{equation}
%     \begin{aligned}
%       p (z)
%         &= (z-\lambda_1) (z-\lambda_2) \cdots (z-\lambda_{m-1}) \underbrace{(z-\lambda_m)(z-\lambda_m)\cdot(z-\lambda_m)}_{n-m+1 \text{ times}} \\
%         &= \left(\sum_{k = 1}^{m-1} (z-\lambda_k) \right) (z-\lambda_m)^{n-m+1} \qquad \forall z \in \myF.
%     \end{aligned}
%   \end{equation}
% \end{xprf}
% If $p$ would have another root than $\lambda_1, \ddd, \lambda_m \in \myF$, say $\lambda_{\text{fresh}}$, we could write $p$ as
% \begin{equation}
%   p(z) = (z- \lambda_{\text{fresh}}) q(z) \quad \forall z \in \myF,
% \end{equation}
% where $q \in \polyn_{n-1} (\myF)$ would have one $(z-\lambda_j)$-term less than $p$. But this is impossible and contradicts the already existing structure of $p$. Hence, $p$ has no other zeros.
